\documentclass[11pt,a4paper]{scrartcl}

\usepackage[T1]{fontenc}
\usepackage[utf8]{inputenc}
\IfFileExists{lithuanian.ldf}{
    \usepackage[lithuanian]{babel}
}{
    \usepackage[provide=*,lithuanian]{babel}
}
% Jei asmeninio stiliaus failas neprieinamas, naudojamas paprastas stilius.
\IfFileExists{../../1.\space Vienanariai\space ir\space daugianariai/manostilius_lt.sty}{
    \usepackage{../../1.\space Vienanariai\space ir\space daugianariai/manostilius_lt}
}{
    \usepackage{amsmath,amssymb,amsthm}
    \usepackage[margin=2.5cm]{geometry}
    \newtheorem{problem}{Uždavinys}
}
\usepackage{tasks}
\renewcommand{\theproblem}{\arabic{problem}}

\title{Pasiruošimas KD}
\author{Andrius Berniukevičius}
\date{}

\begin{document}
\maketitle

\section{Uždaviniai}
Visas lygtis spręskite realiųjų skaičių aibėje. Jei nenurodyta kitaip,
kvadratinių trinarių šaknys ir koeficientai laikomi realiaisiais skaičiais.

\begin{problem}
Išspręskite lygtį:
$$ 4\left(x^2-2x\right)^2+5x^2-10x=-1. $$
\end{problem}

\begin{problem}
Skaičiai $m$ ir $n$ yra daugianario $x^2-2x-10$ šaknys, o $m<n$.
Apskaičiuokite:
\begin{tasks}(2)
    \task $\frac{1}{m}+\frac{1}{n}$;
    \task $n^2-m^2$;
    \task $mn^4+m^4n$;
    \task $\frac{1}{m+1}+\frac{1}{n+1}$.
\end{tasks}
\end{problem}

\begin{problem}
Užrašykite redukuotąją kvadratinę lygtį, kurios sprendiniai yra
$x_1=5\sqrt{2}$ ir $x_2=-\sqrt{18}$.
\end{problem}

\begin{problem}
Išspręskite lygtį:
$$ \frac{x+4}{x^2+x-2}-\frac{3x^2+3x-6}{x+4}=-\frac{1}{2}. $$
\end{problem}

\begin{problem}
Įrodykite, kad kvadratinės lygties $ax^2+bx+c=0$, kur $a\ne0$,
sprendiniai, kai $b^2-4ac\ge0$, apskaičiuojami pagal formules:
$$ x_1=\frac{-b-\sqrt{b^2-4ac}}{2a},\qquad
   x_2=\frac{-b+\sqrt{b^2-4ac}}{2a}. $$
\end{problem}

\begin{problem}
Duotas daugianaris $P(x)=2x^2+4x-1$.
\begin{tasks}(1)
    \task Raskite visas $a$ reikšmes, kurioms
    $(2a+1)\cdot P(a-1)=-3$.
    \task Skaičiai $x_1$ ir $x_2$ yra daugianario $P(x)$ šaknys.
    Apskaičiuokite $4x_1^2+2x_2^2+5x_1+x_2-7$.
    \task Skaičiai $\alpha$ ir $\beta$ yra lygties $P(x)=0$ sprendiniai.
    Apskaičiuokite
    $2\alpha^3+4\beta^3+5\alpha^2+9\beta^2+5\alpha+4\beta-3$.
\end{tasks}
\end{problem}

\begin{problem}
Išskaidykite daugianarį $x^2+6x+5-4(x+1)(x+2)$ dauginamaisiais.
\end{problem}

\begin{problem}
Išspręskite lygtį:
$$ (x^2-5x+1)^2-(x^2+3x+5)^2=0. $$
\end{problem}

\begin{problem}
Duotas kvadratinis trinaris $ax^2+bx+c$, kur $a\ne0$,
o jo šaknys yra $x_1$ ir $x_2$. Įrodykite, kad
$$ x_1+x_2=-\frac{b}{a},\qquad x_1x_2=\frac{c}{a}. $$
\end{problem}

\begin{problem}
Kvadratinio trinario $x^2+bx+18$ šaknų santykis yra $2:1$.
Raskite visas galimas koeficiento $b$ reikšmes.
\end{problem}

\begin{problem}
Išspręskite lygtį:
$$ (x^2+x-6)(x^2+x+1)=-12. $$
\end{problem}

\begin{problem}
Nustatykite, ar yra tokia realioji $a$ reikšmė, kuriai
kvadratinio trinario $P(x)=x^2+(a-1)x+a+2$ realiųjų šaknų
kvadratų suma lygi $9$.
\end{problem}

\begin{problem}
Įrodykite, kad daugianaris $P(x)=x^2+2ax+2a^2+2a+2$
neturi realiųjų šaknų su jokia realiąja $a$ reikšme.
\end{problem}

\begin{problem}
Išspręskite lygtį:
$$ (x^2+x+1)^2+9x^2+6x+1=(6x+2)(x^2+x+1). $$
\end{problem}

\begin{problem}
Išskaidykite reiškinį dauginamaisiais:
$$ 4x^2+12x+9+(2x+3)(3-x)+\frac{16x^4-81}{4x^2+9}. $$
\end{problem}

\begin{problem}
Kvadratinio trinario $-x^2+3x+c$ šaknų skirtumo absoliučioji
vertė lygi $\pi$. Apskaičiuokite koeficiento $c$ reikšmę.
\end{problem}

\begin{problem}
Užrašykite visas redukuotąsias kvadratines lygtis, kurių sprendinių
kvadratų suma lygi $25$, o vienas sprendinys yra vienetu didesnis už kitą.
\end{problem}

\begin{problem}
Suprastinkite trupmeną:
$$ \frac{(3x^2+5x-2)(9x^2-30x+25)}
{(9x^2-18x+5)(10-6x)}. $$
\end{problem}

\begin{problem}
Išspręskite lygtį:
$$ 4^{x+1}-18\cdot2^x+8=0. $$
\end{problem}

\begin{problem}
Kvadratinio trinario $x^2-x-c$ šaknų kubų suma lygi $7$.
Raskite koeficiento $c$ reikšmę.
\end{problem}

\clearpage
\section{Atsakymai}


\begin{enumerate}
    \item $x_1=1-\frac{\sqrt{3}}{2},\quad x_2=1,\quad
    x_3=1+\frac{\sqrt{3}}{2}$.
    \item a) $-\frac15$; b) $4\sqrt{11}$; c) $-680$; d) $-\frac47$.
    \item $x^2-2\sqrt{2}\,x-30=0$.
    \item $x_1=-\frac73,\quad x_2=-\frac32,\quad x_3=0,\quad x_4=2$.
    \item $- $
    \item a) $a_1=-\frac32,\quad a_2=0,\quad a_3=1$;
    b) $2$; c) $-10$.
    \item $-3(x+1)^2$.
    \item $x=-\frac12$.
    \item $-$
    \item $b_1=-9,\quad b_2=9$.
    \item $x_1=\frac{-1-\sqrt{13}}2,\quad x_2=-2,\quad
    x_3=1,\quad x_4=\frac{-1+\sqrt{13}}2$.
    \item $a=-2$.
    \item $P(x)=(x+a)^2+(a+1)^2+1>0$ visiems realiesiems $x$ ir $a$.
    \item $x_1=0,\quad x_2=2$.
    \item $3(2x+3)(x+1)$; reiškinys apibrėžtas visiems $x\in\mathbb R$.
    \item $c=\frac{\pi^2-9}{4}$.
    \item $x^2-7x+12=0$ arba $x^2+7x+12=0$.
    \item $-\frac{x+2}{2}$.
    \item $x_1=-1,\quad x_2=2$.
    \item $c=2$.
\end{enumerate}

\clearpage
\section{Hintai 1}
\begin{enumerate}
    \item Pastebėkite pasikartojantį reiškinį $x^2-2x$.
    \item Taikykite Vijeto teoremą. Skaičių $m$ ir $n$ iš karto
    skaičiuoti nebūtina.
    \item Redukuotojoje kvadratinėje lygtyje koeficientai išreiškiami
    šaknų suma ir sandauga.
    \item Pirmiausia nustatykite apibrėžimo sritį. Atkreipkite dėmesį
    į du vienas kitam atvirkštinius reiškinius.
    \item Abi lygties puses padauginkite iš $4a$.
    \item a) Apskaičiuokite $P(a-1)$.
    b), c) Kiekviena šaknis tenkina lygybę $2t^2+4t-1=0$.
    \item Išskaidykite $x^2+6x+5$ ir ieškokite bendrojo daugiklio.
    \item Taikykite kvadratų skirtumo formulę.
    \item Prisiminkite diskriminanto formulę: kam lygūs
    $x_1$ ir $x_2$?
    \item Šaknis pažymėkite $2t$ ir $t$. Nepamirškite neigiamųjų šaknų.
    \item Pastebėkite abiejuose daugikliuose pasikartojantį reiškinį.
    \item Pagal Vijeto teoremą išreikškite šaknų kvadratų sumą
    per jų sumą ir sandaugą.
    \item Išskirkite pilnąjį kvadratą kintamojo $x$ atžvilgiu.
    \item Pastebėkite, kad $9x^2+6x+1$ yra pilnasis kvadratas.
    \item Atpažinkite pilnąjį kvadratą ir kvadratų skirtumą.
    \item Šaknų skirtumo kvadratą išreikškite jų suma ir sandauga.
    \item Iš duotų sąlygų raskite šaknų sandaugą.
    \item Išskaidykite visus skaitiklio ir vardiklio daugiklius.
    \item Išreikškite $4^{x+1}$ per $2^x$.
    \item Šaknų kubų sumą išreikškite jų suma ir sandauga.
\end{enumerate}

\section{Hintai 2}
\begin{enumerate}
    \item Pažymėkite $t=x^2-2x$. Gausite $4t^2+5t+1=0$.
    \item $m+n=2$, $mn=-10$.
    Naudokite $n^2-m^2=(n-m)(n+m)$ ir
    $mn^4+m^4n=mn(n^3+m^3)$.
    \item Jei šaknys yra $r$ ir $s$, lygtis yra
    $x^2-(r+s)x+rs=0$.
    \item $x\ne-4,-2,1$. Pažymėkite
    $t=\frac{x+4}{x^2+x-2}$; čia $t\ne0$.
    \item Išskirkite pilnąjį kvadratą kairėje pusėje ir perkelkite
    likusius narius į dešinę.
    \item a) $P(a-1)=2a^2-3$.
    b), c) Kiekvienai šakniai $t$ galioja
    $t^2=-2t+\frac12$, o $t^3=\frac92t-1$.
    \item $x^2+6x+5=(x+1)(x+5)$.
    \item Gausite $(-8x-4)(2x^2-2x+6)=0$.
    \item Įstatykite kvadratinės formulės išraiškas į
    $x_1+x_2$ ir $x_1x_2$. Sumai supaprastinkite priešingus
    kvadratinės šaknies narius, o sandaugai pritaikykite
    skirtumo kvadrato formulę.
    \item Pagal Vijeto teoremą $2t^2=18$ ir $3t=-b$.
    \item Pažymėkite $t=x^2+x$. Gausite $(t-6)(t+1)=-12$.
    \item Pagal Vijeto teoremą
    $x_1+x_2=1-a$ ir $x_1x_2=a+2$. Todėl
    \[
    x_1^2+x_2^2=(1-a)^2-2(a+2).
    \]
    Šią išraišką prilyginkite $9$.
    \item $P(x)=(x+a)^2+a^2+2a+2$.
    \item Pažymėkite $A=x^2+x+1$, $B=3x+1$.
    Lygtis yra $A^2+B^2=2AB$.
    \item $4x^2+12x+9=(2x+3)^2$, o
    $16x^4-81=(4x^2-9)(4x^2+9)$.
    \item $x_1+x_2=3$, $x_1x_2=-c$ ir
    $(x_1-x_2)^2=(x_1+x_2)^2-4x_1x_2$.
    \item Kadangi $(x_1-x_2)^2=1$, naudokite tapatybę
    \[
    (x_1-x_2)^2=x_1^2+x_2^2-2x_1x_2
    \]
    ir raskite šaknų sandaugą.
    \item $3x^2+5x-2=(3x-1)(x+2)$,
    $9x^2-18x+5=(3x-1)(3x-5)$.
    \item Pažymėkite $t=2^x>0$. Gausite $4t^2-18t+8=0$.
    \item Jei šaknys yra $r$ ir $s$, tai
    $r^3+s^3=(r+s)^3-3rs(r+s)$.
\end{enumerate}

\section{Hintai 3}
\begin{enumerate}
    \item $(4t+1)(t+1)=0$. Atskirai spręskite
    $x^2-2x=-\frac14$ ir $x^2-2x=-1$.
    \item a) Naudokite $\frac{m+n}{mn}$.
    b) $n-m=\sqrt{(m+n)^2-4mn}$, nes $n>m$.
    c) $m^3+n^3=(m+n)^3-3mn(m+n)$.
    d) Naudokite $\frac{m+n+2}{mn+m+n+1}$.
    \item Šaknų suma lygi $2\sqrt2$, o sandauga lygi $-30$.
    \item $2t^2+t-6=(2t-3)(t+2)=0$.
    Grįžę prie $x$, gausite $3x^2+x-14=0$ arba $2x^2+3x=0$.
    Patikrinkite apibrėžimo sritį.
    \item Kadangi $a\ne0$, lygties abi puses galime padauginti iš
    $4a$:
    \[
    4a^2x^2+4abx+4ac=0.
    \]
    Prie pirmųjų dviejų narių pridėję ir atėmę $b^2$, gauname
    \[
    (2ax+b)^2=b^2-4ac.
    \]
    Kai $b^2-4ac\ge0$, ištraukiame kvadratinę šaknį:
    \[
    2ax+b=\pm\sqrt{b^2-4ac}.
    \]
    Kadangi $2a\ne0$, išreiškiame $x$:
    \[
    x=\frac{-b\pm\sqrt{b^2-4ac}}{2a}.
    \]
    Taigi gauname abi nurodytas sprendinių formules.
    \item a) Lygtis susiveda į $2a(2a+3)(a-1)=0$.
    b) Reiškinys susiveda į $-3(x_1+x_2)-4$.
    c) Reiškinys susiveda į $4(\alpha+\beta)-2$.
    Abiem atvejais šaknų suma lygi $-2$.
    \item Iškelkite $(x+1)$:
    $(x+1)\bigl((x+5)-4(x+2)\bigr)$.
    \item $2x^2-2x+6=2\left(x-\frac12\right)^2+\frac{11}{2}>0$.
    Todėl turi būti $-8x-4=0$.
    \item Pažymėkime $D=b^2-4ac$. Pagal kvadratinę formulę
    \[
    x_1=\frac{-b-\sqrt D}{2a},\qquad
    x_2=\frac{-b+\sqrt D}{2a}.
    \]
    Tada
    \[
    x_1+x_2
    =\frac{-b-\sqrt D-b+\sqrt D}{2a}
    =-\frac ba,
    \]
    o
    \[
    x_1x_2
    =\frac{(-b-\sqrt D)(-b+\sqrt D)}{4a^2}
    =\frac{b^2-D}{4a^2}
    =\frac{4ac}{4a^2}
    =\frac ca.
    \]
    \item $t^2=9$, taigi $t=3$ arba $t=-3$.
    Kiekvieną reikšmę įstatykite į $b=-3t$.
    \item $(t-2)(t-3)=0$. Spręskite $x^2+x=2$ ir $x^2+x=3$.
    \item Iš lygybės
    \[
    (1-a)^2-2(a+2)=9
    \]
    gauname $a^2-4a-12=0$, todėl $a=6$ arba $a=-2$.
    Patikriname, kuri reikšmė duoda realiąsias šaknis:
    \[
    D=(a-1)^2-4(a+2).
    \]
    Kai $a=6$, $D=-7<0$, tad šaknys nėra realiosios.
    Kai $a=-2$, $D=9\ge0$, o trinaris yra
    $P(x)=x^2-3x=x(x-3)$, kurio šaknys $0$ ir $3$.
    Jų kvadratų suma lygi $0^2+3^2=9$, todėl atsakymas $a=-2$.
    \item Užbaikite pertvarkymą:
    $P(x)=(x+a)^2+(a+1)^2+1$. Abu kvadratai neneigiami.
    \item $(A-B)^2=0$, todėl $x^2-2x=0$.
    \item Suprastinta trupmena lygi $4x^2-9=(2x-3)(2x+3)$.
    Iškelkite $(2x+3)$ iš viso reiškinio:
    $(2x+3)\bigl((2x+3)+(3-x)+(2x-3)\bigr)$.
    \item Iš $\pi^2=9+4c$ išreikškite $c$.
    Patikrinkite, kad gautasis trinaris turi realiąsias šaknis.
    \item Tegul šaknys yra $x_1$ ir $x_2$. Iš sąlygų gauname
    \[
    x_1x_2=\frac{25-1}{2}=12,\qquad
    (x_1+x_2)^2=25+2\cdot12=49.
    \]
    Vadinasi, $x_1+x_2=7$ arba $x_1+x_2=-7$. Pagal Vijeto teoremą
    atitinkamos lygtys yra $x^2-7x+12=0$ ir $x^2+7x+12=0$.
    Jų šaknys atitinkamai $3,4$ ir $-3,-4$, tad abiejose lygtyse
    šaknų kvadratų suma yra $25$, o skirtumo absoliučioji vertė — $1$.
    \item $9x^2-30x+25=(3x-5)^2$ ir $10-6x=-2(3x-5)$.
    Sutraukite bendruosius daugiklius.
    \item $4t^2-18t+8=2(2t-1)(t-4)$.
    Spręskite $2^x=\frac12$ ir $2^x=4$.
    \item Pagal Vijeto teoremą $r+s=1$, $rs=-c$.
    Todėl kubų sumos sąlyga virsta $1+3c=7$.
\end{enumerate}
\end{document}